\documentclass[11pt,a4paper,reqno]{amsart}
\usepackage[T1]{fontenc}
\usepackage{lmodern}
\usepackage[margin=25mm]{geometry}
\usepackage{amsmath,amssymb,amsthm,listings,booktabs,array,tikz}
\usepackage[expansion=false]{microtype}
\setlength{\emergencystretch}{1em}
\usepackage[hidelinks]{hyperref}
\hypersetup{pdftitle={A new lower bound for the sixth Schur number: S(6) >= 537},pdfauthor={Zhiyang Chen, Ligeng Zhu, Fan Long}}
\lstset{basicstyle=\small\ttfamily,columns=fullflexible,keepspaces=true,
  breaklines=true,showstringspaces=false,aboveskip=4pt,belowskip=4pt}
\newtheorem{theorem}{Theorem}
\newtheorem{proposition}[theorem]{Proposition}
\newtheorem{corollary}[theorem]{Corollary}
\theoremstyle{definition}
\newtheorem{definition}[theorem]{Definition}
\newtheorem{remark}[theorem]{Remark}
\newcommand{\fold}{\operatorname{f}}

\title[The sixth Schur number]{A new lower bound for the sixth Schur number: $S(6)\ge537$}
\author{Zhiyang Chen}
\author{Ligeng Zhu}
\author{Fan Long}
\date{22 October 2026}

\begin{document}

\begin{abstract}
We give an explicit partition of $\{1,\ldots,537\}$ into six sum-free sets,
establishing $S(6)\ge537$. This improves the lower bound 536 of Fredricksen and
Sweet (2000). Known recursive constructions then give improved lower bounds
for $S(8)$, $S(11)$, $S(13)$, and $R_6(3)$. We also report 131 pairwise
inequivalent partitions found by the search. We describe their structure and
provide the complete certificates and verification code in a companion repository.
\end{abstract}
\keywords{Schur numbers, sum-free partitions, SAT, computational combinatorics}
\maketitle

\section{Introduction}
A set $C$ of integers is \emph{sum-free} if there are no $x,y\in C$, not necessarily
distinct, with $x+y\in C$. The Schur number $S(k)$ is the largest $n$ for which
$\{1,\dots,n\}$ can be partitioned into $k$ sum-free sets. The values
$S(1),\dots,S(5)=1,4,13,44,160$ are known, the last one by a large SAT computation of
Heule~\cite{heule}. For $k=6$ the record has been $S(6)\ge536$ since 2000, when
Fredricksen and Sweet~\cite{fs} found it by a search over symmetric sum-free
partitions; it has been quoted as the best known bound ever since. We improve it
by one.

\begin{theorem}\label{thm:537}
$S(6)\ge 537$.
\end{theorem}

\begin{proof}
The following six sets are pairwise disjoint, their union is $\{1,\dots,537\}$, and
none of them contains $x$, $y$ and $x+y$. We list every element, as in the
explicit constructions of Fredricksen and Sweet~\cite{fs}. This is a finite check; a script that
performs it is supplied as \texttt{check.py} in the companion repository.
\par\smallskip
\noindent\begin{minipage}{\linewidth}
\textbf{Color 1 (76 integers).}
\begin{lstlisting}
  1   5  12  15  19  21  23  25  29  32  39  43  45  49  95
102 106 108 112 119 122 126 128 130 132 136 139 146 150 152
156 163 166 170 172 174 176 180 358 362 364 366 368 372 375
382 386 388 392 399 402 406 408 410 412 416 419 426 430 432
436 443 489 493 495 499 506 509 513 515 517 519 523 526 533
537
\end{lstlisting}
\end{minipage}\par\smallskip

\noindent\begin{minipage}{\linewidth}
\textbf{Color 2 (76 integers).}
\begin{lstlisting}
  2   6   7  17  18  22  27  37  41  42  46  51  66  70  81
 89  90 100 101 105 109 110 113 114 124 125 129 133 134 144
145 149 153 157 158 168 169 192 346 369 370 380 381 385 389
393 394 404 405 409 413 414 424 425 428 429 433 437 438 448
449 457 468 472 487 492 496 497 501 511 516 520 521 531 532
536
\end{lstlisting}
\end{minipage}\par\smallskip

\noindent\begin{minipage}{\linewidth}
\textbf{Color 3 (88 integers).}
\begin{lstlisting}
  3   8   9  14  20  24  26  30  31  36  47  53  59  63  65
 69  76  82  86  88  92  98 104 115 120 121 127 131 137 142
143 148 154 159 160 165 171 177 181 182 187 188 193 204 334
345 350 351 356 357 361 367 373 378 379 384 390 395 396 401
407 411 417 418 423 434 440 446 450 452 456 462 469 473 475
479 485 491 502 507 508 512 514 518 524 529 530 535
\end{lstlisting}
\end{minipage}\par\smallskip

\noindent\begin{minipage}{\linewidth}
\textbf{Color 4 (70 integers).}
\begin{lstlisting}
  4  10  11  13  16  28  33  34  35  40  55  57  79  93  94
 96  99 111 116 117 118 123 135 138 140 141 147 155 161 162
164 167 179 185 203 335 353 359 371 374 376 377 383 391 397
398 400 403 415 420 421 422 427 439 442 444 445 459 481 483
498 503 504 505 510 522 525 527 528 534
\end{lstlisting}
\end{minipage}\par\smallskip

\noindent\begin{minipage}{\linewidth}
\textbf{Color 5 (76 integers).}
\begin{lstlisting}
 38  44  48  50  52  54  56  58  60  61  62  64  67  68  71
 72  73  74  75  77  78  80  83  84  85  87  91  97 103 107
173 189 196 197 201 209 212 236 302 326 329 337 341 342 349
365 431 435 441 447 451 453 454 455 458 460 461 463 464 465
466 467 470 471 474 476 477 478 480 482 484 486 488 490 494
500
\end{lstlisting}
\end{minipage}\par\smallskip

\noindent\begin{minipage}{\linewidth}
\textbf{Color 6 (151 integers).}
\begin{lstlisting}
151 175 178 183 184 186 190 191 194 195 198 199 200 202 205
206 207 208 210 211 213 214 215 216 217 218 219 220 221 222
223 224 225 226 227 228 229 230 231 232 233 234 235 237 238
239 240 241 242 243 244 245 246 247 248 249 250 251 252 253
254 255 256 257 258 259 260 261 262 263 264 265 266 267 268
269 270 271 272 273 274 275 276 277 278 279 280 281 282 283
284 285 286 287 288 289 290 291 292 293 294 295 296 297 298
299 300 301 303 304 305 306 307 308 309 310 311 312 313 314
315 316 317 318 319 320 321 322 323 324 325 327 328 330 331
332 333 336 338 339 340 343 344 347 348 352 354 355 360 363
387
\end{lstlisting}
\end{minipage}\par\smallskip

\end{proof}

Theorem~\ref{thm:537} needs only this one partition. Section~\ref{sec:consequences}
lists what it implies for other Schur and Ramsey numbers. In Section~\ref{sec:families} we
describe all $131$ partitions of $\{1,\dots,537\}$ into six sum-free sets that our
search produced.
Section~\ref{sec:method} explains briefly how the partitions were found.
The complete results, code, and supplementary tables are available at
\url{https://github.com/jeffchen006/Schur-Number-6}.

Throughout, $N=538$, $\fold(x)=\min(x,N-x)$, and a colouring $c$ of
$\{1,\dots,537\}$ is \emph{palindromic} if $c(x)=c(N-x)$ for all $x$. We number the
colours $1,\dots,6$ in the order of their least elements.

\section{Consequences for other bounds}\label{sec:consequences}
Classical recurrences of Schur and Abbott--Hanson~\cite{schur,abbott},
and the template constructions of Rowley and Ageron et al.~\cite{rowley,ageron},
turn smaller partitions into larger ones. We use
\[
S(k+2)\ge10S(k)+2\quad\text{\cite{bengone}},\qquad
S(k+5)\ge380S(k)+148\quad\text{\cite{ageron}}.
\]
A sum-free $k$-partition of $[1,n]$ also gives an edge coloring of $K_{n+1}$
without monochromatic triangles: color $ij$ by the class of $|i-j|$.
Consequently, $R_k(3)\ge S(k)+2$.

\begin{corollary}\label{cor:consequences}
The lower bounds in the column headed ``This work'' in Table~\ref{tab:bounds} hold.
\end{corollary}

\begin{table}[tbp]
\caption{Lower bounds in the cited literature and the bounds obtained here.
The six- and eight-color bounds have explicit certificates in the repository.}
\label{tab:bounds}
\centering
\begin{tabular}{lrrl}
\toprule
Quantity & Previous bound & This work & Previous source\\
\midrule
$S(6)$    & $536$          & $537$            & Fredricksen--Sweet~\cite{fs}\\
$S(8)$    & $5\,362$       & $5\,372$         & Bengone et al.~\cite{bengone}\\
$S(11)$   & $203\,828$     & $204\,208$       & Ageron et al.~\cite{ageron}\\
$S(13)$   & $2\,038\,282$ & $2\,042\,082$    & Bengone et al.~\cite{bengone}\\
$R_6(3)$  & $538$          & $539$            & Fredricksen--Sweet~\cite{fs,xxer}\\
\bottomrule
\end{tabular}
\end{table}

\begin{proof}
$S(8)\ge10\cdot537+2$, $S(11)\ge380\cdot537+148$,
$S(13)\ge10\cdot S(11)+2\ge10\cdot204\,208+2$, and $R_6(3)\ge S(6)+2$. For $S(8)$ we also
built the partition explicitly: apply the width-$10$ template of
\cite[Table~II]{bengone} to the partition of Theorem~\ref{thm:537}, with colours $7$
and $8$ for the labels $A$ and $B$. The companion script \texttt{check.py} builds this
partition of $\{1,\dots,5372\}$ into eight sets and checks that all of them are
sum-free.
\end{proof}

Bounds for larger $k$ that are derived through these recurrences change in the same
way; we list only the four cases in which we checked the previously published
values.

\section{The 537-partitions and their families}\label{sec:families}
The example above is palindromic, and each of its classes is even sum-free modulo
$538$. The least elements of its classes are $1,2,3,4,38,151$, and the sixth class
contains the whole interval $[237,301]$. These features are captured by the
following family of colourings.

\begin{definition}\label{def:frameg}
A colouring $c\colon\{1,\dots,537\}\to\{1,\dots,6\}$ whose classes are sum-free
modulo $538$ is in \emph{frame $g$} if, with $Q=151$ and $d=38$,
\begin{align*}
 \fold(x)<j &\ \Rightarrow\ c(x)\neq j \quad (j=1,2,3,4),\\
 \fold(x)<d \ \text{or}\ Q-d<\fold(x)<Q &\ \Rightarrow\ c(x)\neq 5,\\
 \fold(x)\in\{d,\,N-2Q\} &\ \Rightarrow\ c(x)=5,\\
 \fold(x)<Q &\ \Rightarrow\ c(x)\neq 6,\\
 \fold(x)=Q \ \text{or}\ \fold(x)>N-2Q &\ \Rightarrow\ c(x)=6.
\end{align*}
\end{definition}

In words: the sixth colour does not occur below $151$ and fills the middle of the
interval, and the fifth colour does not occur below $38$ or just below $151$.
Figure~\ref{fig:structure} shows the certificate and these restrictions.

\begin{figure}[tbp]
\centering
% Generated by make_figures.py from certificate.col; do not edit data by hand.
\definecolor{schur1}{HTML}{0072B2}
\definecolor{schur2}{HTML}{E69F00}
\definecolor{schur3}{HTML}{009E73}
\definecolor{schur4}{HTML}{CC79A7}
\definecolor{schur5}{HTML}{D55E00}
\definecolor{schur6}{HTML}{333333}
\begin{tikzpicture}[x=1cm,y=1cm,font=\scriptsize]
\node[anchor=west,font=\small] at (-0.65,0.43) {(a) The explicit coloring on $[1,537]$};
\node[anchor=east] at (-0.12,0.0900) {$C_1$};
\fill[black!4] (0,0.0000) rectangle (13.1,0.1800);
\fill[schur1] (0.000000,0.0000) rectangle (0.024395,0.1800);
\fill[schur1] (0.097579,0.0000) rectangle (0.121974,0.1800);
\fill[schur1] (0.268343,0.0000) rectangle (0.292737,0.1800);
\fill[schur1] (0.341527,0.0000) rectangle (0.365922,0.1800);
\fill[schur1] (0.439106,0.0000) rectangle (0.463501,0.1800);
\fill[schur1] (0.487896,0.0000) rectangle (0.512291,0.1800);
\fill[schur1] (0.536685,0.0000) rectangle (0.561080,0.1800);
\fill[schur1] (0.585475,0.0000) rectangle (0.609870,0.1800);
\fill[schur1] (0.683054,0.0000) rectangle (0.707449,0.1800);
\fill[schur1] (0.756238,0.0000) rectangle (0.780633,0.1800);
\fill[schur1] (0.927002,0.0000) rectangle (0.951397,0.1800);
\fill[schur1] (1.024581,0.0000) rectangle (1.048976,0.1800);
\fill[schur1] (1.073371,0.0000) rectangle (1.097765,0.1800);
\fill[schur1] (1.170950,0.0000) rectangle (1.195345,0.1800);
\fill[schur1] (2.293110,0.0000) rectangle (2.317505,0.1800);
\fill[schur1] (2.463873,0.0000) rectangle (2.488268,0.1800);
\fill[schur1] (2.561453,0.0000) rectangle (2.585847,0.1800);
\fill[schur1] (2.610242,0.0000) rectangle (2.634637,0.1800);
\fill[schur1] (2.707821,0.0000) rectangle (2.732216,0.1800);
\fill[schur1] (2.878585,0.0000) rectangle (2.902980,0.1800);
\fill[schur1] (2.951769,0.0000) rectangle (2.976164,0.1800);
\fill[schur1] (3.049348,0.0000) rectangle (3.073743,0.1800);
\fill[schur1] (3.098138,0.0000) rectangle (3.122533,0.1800);
\fill[schur1] (3.146927,0.0000) rectangle (3.171322,0.1800);
\fill[schur1] (3.195717,0.0000) rectangle (3.220112,0.1800);
\fill[schur1] (3.293296,0.0000) rectangle (3.317691,0.1800);
\fill[schur1] (3.366480,0.0000) rectangle (3.390875,0.1800);
\fill[schur1] (3.537244,0.0000) rectangle (3.561639,0.1800);
\fill[schur1] (3.634823,0.0000) rectangle (3.659218,0.1800);
\fill[schur1] (3.683613,0.0000) rectangle (3.708007,0.1800);
\fill[schur1] (3.781192,0.0000) rectangle (3.805587,0.1800);
\fill[schur1] (3.951955,0.0000) rectangle (3.976350,0.1800);
\fill[schur1] (4.025140,0.0000) rectangle (4.049534,0.1800);
\fill[schur1] (4.122719,0.0000) rectangle (4.147114,0.1800);
\fill[schur1] (4.171508,0.0000) rectangle (4.195903,0.1800);
\fill[schur1] (4.220298,0.0000) rectangle (4.244693,0.1800);
\fill[schur1] (4.269088,0.0000) rectangle (4.293482,0.1800);
\fill[schur1] (4.366667,0.0000) rectangle (4.391061,0.1800);
\fill[schur1] (8.708939,0.0000) rectangle (8.733333,0.1800);
\fill[schur1] (8.806518,0.0000) rectangle (8.830912,0.1800);
\fill[schur1] (8.855307,0.0000) rectangle (8.879702,0.1800);
\fill[schur1] (8.904097,0.0000) rectangle (8.928492,0.1800);
\fill[schur1] (8.952886,0.0000) rectangle (8.977281,0.1800);
\fill[schur1] (9.050466,0.0000) rectangle (9.074860,0.1800);
\fill[schur1] (9.123650,0.0000) rectangle (9.148045,0.1800);
\fill[schur1] (9.294413,0.0000) rectangle (9.318808,0.1800);
\fill[schur1] (9.391993,0.0000) rectangle (9.416387,0.1800);
\fill[schur1] (9.440782,0.0000) rectangle (9.465177,0.1800);
\fill[schur1] (9.538361,0.0000) rectangle (9.562756,0.1800);
\fill[schur1] (9.709125,0.0000) rectangle (9.733520,0.1800);
\fill[schur1] (9.782309,0.0000) rectangle (9.806704,0.1800);
\fill[schur1] (9.879888,0.0000) rectangle (9.904283,0.1800);
\fill[schur1] (9.928678,0.0000) rectangle (9.953073,0.1800);
\fill[schur1] (9.977467,0.0000) rectangle (10.001862,0.1800);
\fill[schur1] (10.026257,0.0000) rectangle (10.050652,0.1800);
\fill[schur1] (10.123836,0.0000) rectangle (10.148231,0.1800);
\fill[schur1] (10.197020,0.0000) rectangle (10.221415,0.1800);
\fill[schur1] (10.367784,0.0000) rectangle (10.392179,0.1800);
\fill[schur1] (10.465363,0.0000) rectangle (10.489758,0.1800);
\fill[schur1] (10.514153,0.0000) rectangle (10.538547,0.1800);
\fill[schur1] (10.611732,0.0000) rectangle (10.636127,0.1800);
\fill[schur1] (10.782495,0.0000) rectangle (10.806890,0.1800);
\fill[schur1] (11.904655,0.0000) rectangle (11.929050,0.1800);
\fill[schur1] (12.002235,0.0000) rectangle (12.026629,0.1800);
\fill[schur1] (12.051024,0.0000) rectangle (12.075419,0.1800);
\fill[schur1] (12.148603,0.0000) rectangle (12.172998,0.1800);
\fill[schur1] (12.319367,0.0000) rectangle (12.343762,0.1800);
\fill[schur1] (12.392551,0.0000) rectangle (12.416946,0.1800);
\fill[schur1] (12.490130,0.0000) rectangle (12.514525,0.1800);
\fill[schur1] (12.538920,0.0000) rectangle (12.563315,0.1800);
\fill[schur1] (12.587709,0.0000) rectangle (12.612104,0.1800);
\fill[schur1] (12.636499,0.0000) rectangle (12.660894,0.1800);
\fill[schur1] (12.734078,0.0000) rectangle (12.758473,0.1800);
\fill[schur1] (12.807263,0.0000) rectangle (12.831657,0.1800);
\fill[schur1] (12.978026,0.0000) rectangle (13.002421,0.1800);
\fill[schur1] (13.075605,0.0000) rectangle (13.100000,0.1800);
\node[anchor=west] at (13.219999999999999,0.0900) {76};
\node[anchor=east] at (-0.12,-0.1800) {$C_2$};
\fill[black!4] (0,-0.2700) rectangle (13.1,-0.0900);
\fill[schur2] (0.024395,-0.2700) rectangle (0.048790,-0.0900);
\fill[schur2] (0.121974,-0.2700) rectangle (0.170764,-0.0900);
\fill[schur2] (0.390317,-0.2700) rectangle (0.439106,-0.0900);
\fill[schur2] (0.512291,-0.2700) rectangle (0.536685,-0.0900);
\fill[schur2] (0.634264,-0.2700) rectangle (0.658659,-0.0900);
\fill[schur2] (0.878212,-0.2700) rectangle (0.902607,-0.0900);
\fill[schur2] (0.975791,-0.2700) rectangle (1.024581,-0.0900);
\fill[schur2] (1.097765,-0.2700) rectangle (1.122160,-0.0900);
\fill[schur2] (1.219739,-0.2700) rectangle (1.244134,-0.0900);
\fill[schur2] (1.585661,-0.2700) rectangle (1.610056,-0.0900);
\fill[schur2] (1.683240,-0.2700) rectangle (1.707635,-0.0900);
\fill[schur2] (1.951583,-0.2700) rectangle (1.975978,-0.0900);
\fill[schur2] (2.146741,-0.2700) rectangle (2.195531,-0.0900);
\fill[schur2] (2.415084,-0.2700) rectangle (2.463873,-0.0900);
\fill[schur2] (2.537058,-0.2700) rectangle (2.561453,-0.0900);
\fill[schur2] (2.634637,-0.2700) rectangle (2.683426,-0.0900);
\fill[schur2] (2.732216,-0.2700) rectangle (2.781006,-0.0900);
\fill[schur2] (3.000559,-0.2700) rectangle (3.049348,-0.0900);
\fill[schur2] (3.122533,-0.2700) rectangle (3.146927,-0.0900);
\fill[schur2] (3.220112,-0.2700) rectangle (3.268901,-0.0900);
\fill[schur2] (3.488454,-0.2700) rectangle (3.537244,-0.0900);
\fill[schur2] (3.610428,-0.2700) rectangle (3.634823,-0.0900);
\fill[schur2] (3.708007,-0.2700) rectangle (3.732402,-0.0900);
\fill[schur2] (3.805587,-0.2700) rectangle (3.854376,-0.0900);
\fill[schur2] (4.073929,-0.2700) rectangle (4.122719,-0.0900);
\fill[schur2] (4.659404,-0.2700) rectangle (4.683799,-0.0900);
\fill[schur2] (8.416201,-0.2700) rectangle (8.440596,-0.0900);
\fill[schur2] (8.977281,-0.2700) rectangle (9.026071,-0.0900);
\fill[schur2] (9.245624,-0.2700) rectangle (9.294413,-0.0900);
\fill[schur2] (9.367598,-0.2700) rectangle (9.391993,-0.0900);
\fill[schur2] (9.465177,-0.2700) rectangle (9.489572,-0.0900);
\fill[schur2] (9.562756,-0.2700) rectangle (9.611546,-0.0900);
\fill[schur2] (9.831099,-0.2700) rectangle (9.879888,-0.0900);
\fill[schur2] (9.953073,-0.2700) rectangle (9.977467,-0.0900);
\fill[schur2] (10.050652,-0.2700) rectangle (10.099441,-0.0900);
\fill[schur2] (10.318994,-0.2700) rectangle (10.367784,-0.0900);
\fill[schur2] (10.416574,-0.2700) rectangle (10.465363,-0.0900);
\fill[schur2] (10.538547,-0.2700) rectangle (10.562942,-0.0900);
\fill[schur2] (10.636127,-0.2700) rectangle (10.684916,-0.0900);
\fill[schur2] (10.904469,-0.2700) rectangle (10.953259,-0.0900);
\fill[schur2] (11.124022,-0.2700) rectangle (11.148417,-0.0900);
\fill[schur2] (11.392365,-0.2700) rectangle (11.416760,-0.0900);
\fill[schur2] (11.489944,-0.2700) rectangle (11.514339,-0.0900);
\fill[schur2] (11.855866,-0.2700) rectangle (11.880261,-0.0900);
\fill[schur2] (11.977840,-0.2700) rectangle (12.002235,-0.0900);
\fill[schur2] (12.075419,-0.2700) rectangle (12.124209,-0.0900);
\fill[schur2] (12.197393,-0.2700) rectangle (12.221788,-0.0900);
\fill[schur2] (12.441341,-0.2700) rectangle (12.465736,-0.0900);
\fill[schur2] (12.563315,-0.2700) rectangle (12.587709,-0.0900);
\fill[schur2] (12.660894,-0.2700) rectangle (12.709683,-0.0900);
\fill[schur2] (12.929236,-0.2700) rectangle (12.978026,-0.0900);
\fill[schur2] (13.051210,-0.2700) rectangle (13.075605,-0.0900);
\node[anchor=west] at (13.219999999999999,-0.1800) {76};
\node[anchor=east] at (-0.12,-0.4500) {$C_3$};
\fill[black!4] (0,-0.5400) rectangle (13.1,-0.3600);
\fill[schur3] (0.048790,-0.5400) rectangle (0.073184,-0.3600);
\fill[schur3] (0.170764,-0.5400) rectangle (0.219553,-0.3600);
\fill[schur3] (0.317132,-0.5400) rectangle (0.341527,-0.3600);
\fill[schur3] (0.463501,-0.5400) rectangle (0.487896,-0.3600);
\fill[schur3] (0.561080,-0.5400) rectangle (0.585475,-0.3600);
\fill[schur3] (0.609870,-0.5400) rectangle (0.634264,-0.3600);
\fill[schur3] (0.707449,-0.5400) rectangle (0.756238,-0.3600);
\fill[schur3] (0.853818,-0.5400) rectangle (0.878212,-0.3600);
\fill[schur3] (1.122160,-0.5400) rectangle (1.146555,-0.3600);
\fill[schur3] (1.268529,-0.5400) rectangle (1.292924,-0.3600);
\fill[schur3] (1.414898,-0.5400) rectangle (1.439292,-0.3600);
\fill[schur3] (1.512477,-0.5400) rectangle (1.536872,-0.3600);
\fill[schur3] (1.561266,-0.5400) rectangle (1.585661,-0.3600);
\fill[schur3] (1.658845,-0.5400) rectangle (1.683240,-0.3600);
\fill[schur3] (1.829609,-0.5400) rectangle (1.854004,-0.3600);
\fill[schur3] (1.975978,-0.5400) rectangle (2.000372,-0.3600);
\fill[schur3] (2.073557,-0.5400) rectangle (2.097952,-0.3600);
\fill[schur3] (2.122346,-0.5400) rectangle (2.146741,-0.3600);
\fill[schur3] (2.219926,-0.5400) rectangle (2.244320,-0.3600);
\fill[schur3] (2.366294,-0.5400) rectangle (2.390689,-0.3600);
\fill[schur3] (2.512663,-0.5400) rectangle (2.537058,-0.3600);
\fill[schur3] (2.781006,-0.5400) rectangle (2.805400,-0.3600);
\fill[schur3] (2.902980,-0.5400) rectangle (2.951769,-0.3600);
\fill[schur3] (3.073743,-0.5400) rectangle (3.098138,-0.3600);
\fill[schur3] (3.171322,-0.5400) rectangle (3.195717,-0.3600);
\fill[schur3] (3.317691,-0.5400) rectangle (3.342086,-0.3600);
\fill[schur3] (3.439665,-0.5400) rectangle (3.488454,-0.3600);
\fill[schur3] (3.586034,-0.5400) rectangle (3.610428,-0.3600);
\fill[schur3] (3.732402,-0.5400) rectangle (3.756797,-0.3600);
\fill[schur3] (3.854376,-0.5400) rectangle (3.903166,-0.3600);
\fill[schur3] (4.000745,-0.5400) rectangle (4.025140,-0.3600);
\fill[schur3] (4.147114,-0.5400) rectangle (4.171508,-0.3600);
\fill[schur3] (4.293482,-0.5400) rectangle (4.317877,-0.3600);
\fill[schur3] (4.391061,-0.5400) rectangle (4.439851,-0.3600);
\fill[schur3] (4.537430,-0.5400) rectangle (4.586220,-0.3600);
\fill[schur3] (4.683799,-0.5400) rectangle (4.708194,-0.3600);
\fill[schur3] (4.952142,-0.5400) rectangle (4.976536,-0.3600);
\fill[schur3] (8.123464,-0.5400) rectangle (8.147858,-0.3600);
\fill[schur3] (8.391806,-0.5400) rectangle (8.416201,-0.3600);
\fill[schur3] (8.513780,-0.5400) rectangle (8.562570,-0.3600);
\fill[schur3] (8.660149,-0.5400) rectangle (8.708939,-0.3600);
\fill[schur3] (8.782123,-0.5400) rectangle (8.806518,-0.3600);
\fill[schur3] (8.928492,-0.5400) rectangle (8.952886,-0.3600);
\fill[schur3] (9.074860,-0.5400) rectangle (9.099255,-0.3600);
\fill[schur3] (9.196834,-0.5400) rectangle (9.245624,-0.3600);
\fill[schur3] (9.343203,-0.5400) rectangle (9.367598,-0.3600);
\fill[schur3] (9.489572,-0.5400) rectangle (9.513966,-0.3600);
\fill[schur3] (9.611546,-0.5400) rectangle (9.660335,-0.3600);
\fill[schur3] (9.757914,-0.5400) rectangle (9.782309,-0.3600);
\fill[schur3] (9.904283,-0.5400) rectangle (9.928678,-0.3600);
\fill[schur3] (10.001862,-0.5400) rectangle (10.026257,-0.3600);
\fill[schur3] (10.148231,-0.5400) rectangle (10.197020,-0.3600);
\fill[schur3] (10.294600,-0.5400) rectangle (10.318994,-0.3600);
\fill[schur3] (10.562942,-0.5400) rectangle (10.587337,-0.3600);
\fill[schur3] (10.709311,-0.5400) rectangle (10.733706,-0.3600);
\fill[schur3] (10.855680,-0.5400) rectangle (10.880074,-0.3600);
\fill[schur3] (10.953259,-0.5400) rectangle (10.977654,-0.3600);
\fill[schur3] (11.002048,-0.5400) rectangle (11.026443,-0.3600);
\fill[schur3] (11.099628,-0.5400) rectangle (11.124022,-0.3600);
\fill[schur3] (11.245996,-0.5400) rectangle (11.270391,-0.3600);
\fill[schur3] (11.416760,-0.5400) rectangle (11.441155,-0.3600);
\fill[schur3] (11.514339,-0.5400) rectangle (11.538734,-0.3600);
\fill[schur3] (11.563128,-0.5400) rectangle (11.587523,-0.3600);
\fill[schur3] (11.660708,-0.5400) rectangle (11.685102,-0.3600);
\fill[schur3] (11.807076,-0.5400) rectangle (11.831471,-0.3600);
\fill[schur3] (11.953445,-0.5400) rectangle (11.977840,-0.3600);
\fill[schur3] (12.221788,-0.5400) rectangle (12.246182,-0.3600);
\fill[schur3] (12.343762,-0.5400) rectangle (12.392551,-0.3600);
\fill[schur3] (12.465736,-0.5400) rectangle (12.490130,-0.3600);
\fill[schur3] (12.514525,-0.5400) rectangle (12.538920,-0.3600);
\fill[schur3] (12.612104,-0.5400) rectangle (12.636499,-0.3600);
\fill[schur3] (12.758473,-0.5400) rectangle (12.782868,-0.3600);
\fill[schur3] (12.880447,-0.5400) rectangle (12.929236,-0.3600);
\fill[schur3] (13.026816,-0.5400) rectangle (13.051210,-0.3600);
\node[anchor=west] at (13.219999999999999,-0.4500) {88};
\node[anchor=east] at (-0.12,-0.7200) {$C_4$};
\fill[black!4] (0,-0.8100) rectangle (13.1,-0.6300);
\fill[schur4] (0.073184,-0.8100) rectangle (0.097579,-0.6300);
\fill[schur4] (0.219553,-0.8100) rectangle (0.268343,-0.6300);
\fill[schur4] (0.292737,-0.8100) rectangle (0.317132,-0.6300);
\fill[schur4] (0.365922,-0.8100) rectangle (0.390317,-0.6300);
\fill[schur4] (0.658659,-0.8100) rectangle (0.683054,-0.6300);
\fill[schur4] (0.780633,-0.8100) rectangle (0.853818,-0.6300);
\fill[schur4] (0.951397,-0.8100) rectangle (0.975791,-0.6300);
\fill[schur4] (1.317318,-0.8100) rectangle (1.341713,-0.6300);
\fill[schur4] (1.366108,-0.8100) rectangle (1.390503,-0.6300);
\fill[schur4] (1.902793,-0.8100) rectangle (1.927188,-0.6300);
\fill[schur4] (2.244320,-0.8100) rectangle (2.293110,-0.6300);
\fill[schur4] (2.317505,-0.8100) rectangle (2.341899,-0.6300);
\fill[schur4] (2.390689,-0.8100) rectangle (2.415084,-0.6300);
\fill[schur4] (2.683426,-0.8100) rectangle (2.707821,-0.6300);
\fill[schur4] (2.805400,-0.8100) rectangle (2.878585,-0.6300);
\fill[schur4] (2.976164,-0.8100) rectangle (3.000559,-0.6300);
\fill[schur4] (3.268901,-0.8100) rectangle (3.293296,-0.6300);
\fill[schur4] (3.342086,-0.8100) rectangle (3.366480,-0.6300);
\fill[schur4] (3.390875,-0.8100) rectangle (3.439665,-0.6300);
\fill[schur4] (3.561639,-0.8100) rectangle (3.586034,-0.6300);
\fill[schur4] (3.756797,-0.8100) rectangle (3.781192,-0.6300);
\fill[schur4] (3.903166,-0.8100) rectangle (3.951955,-0.6300);
\fill[schur4] (3.976350,-0.8100) rectangle (4.000745,-0.6300);
\fill[schur4] (4.049534,-0.8100) rectangle (4.073929,-0.6300);
\fill[schur4] (4.342272,-0.8100) rectangle (4.366667,-0.6300);
\fill[schur4] (4.488641,-0.8100) rectangle (4.513035,-0.6300);
\fill[schur4] (4.927747,-0.8100) rectangle (4.952142,-0.6300);
\fill[schur4] (8.147858,-0.8100) rectangle (8.172253,-0.6300);
\fill[schur4] (8.586965,-0.8100) rectangle (8.611359,-0.6300);
\fill[schur4] (8.733333,-0.8100) rectangle (8.757728,-0.6300);
\fill[schur4] (9.026071,-0.8100) rectangle (9.050466,-0.6300);
\fill[schur4] (9.099255,-0.8100) rectangle (9.123650,-0.6300);
\fill[schur4] (9.148045,-0.8100) rectangle (9.196834,-0.6300);
\fill[schur4] (9.318808,-0.8100) rectangle (9.343203,-0.6300);
\fill[schur4] (9.513966,-0.8100) rectangle (9.538361,-0.6300);
\fill[schur4] (9.660335,-0.8100) rectangle (9.709125,-0.6300);
\fill[schur4] (9.733520,-0.8100) rectangle (9.757914,-0.6300);
\fill[schur4] (9.806704,-0.8100) rectangle (9.831099,-0.6300);
\fill[schur4] (10.099441,-0.8100) rectangle (10.123836,-0.6300);
\fill[schur4] (10.221415,-0.8100) rectangle (10.294600,-0.6300);
\fill[schur4] (10.392179,-0.8100) rectangle (10.416574,-0.6300);
\fill[schur4] (10.684916,-0.8100) rectangle (10.709311,-0.6300);
\fill[schur4] (10.758101,-0.8100) rectangle (10.782495,-0.6300);
\fill[schur4] (10.806890,-0.8100) rectangle (10.855680,-0.6300);
\fill[schur4] (11.172812,-0.8100) rectangle (11.197207,-0.6300);
\fill[schur4] (11.709497,-0.8100) rectangle (11.733892,-0.6300);
\fill[schur4] (11.758287,-0.8100) rectangle (11.782682,-0.6300);
\fill[schur4] (12.124209,-0.8100) rectangle (12.148603,-0.6300);
\fill[schur4] (12.246182,-0.8100) rectangle (12.319367,-0.6300);
\fill[schur4] (12.416946,-0.8100) rectangle (12.441341,-0.6300);
\fill[schur4] (12.709683,-0.8100) rectangle (12.734078,-0.6300);
\fill[schur4] (12.782868,-0.8100) rectangle (12.807263,-0.6300);
\fill[schur4] (12.831657,-0.8100) rectangle (12.880447,-0.6300);
\fill[schur4] (13.002421,-0.8100) rectangle (13.026816,-0.6300);
\node[anchor=west] at (13.219999999999999,-0.7200) {70};
\node[anchor=east] at (-0.12,-0.9900) {$C_5$};
\fill[black!4] (0,-1.0800) rectangle (13.1,-0.9000);
\fill[schur5] (0.902607,-1.0800) rectangle (0.927002,-0.9000);
\fill[schur5] (1.048976,-1.0800) rectangle (1.073371,-0.9000);
\fill[schur5] (1.146555,-1.0800) rectangle (1.170950,-0.9000);
\fill[schur5] (1.195345,-1.0800) rectangle (1.219739,-0.9000);
\fill[schur5] (1.244134,-1.0800) rectangle (1.268529,-0.9000);
\fill[schur5] (1.292924,-1.0800) rectangle (1.317318,-0.9000);
\fill[schur5] (1.341713,-1.0800) rectangle (1.366108,-0.9000);
\fill[schur5] (1.390503,-1.0800) rectangle (1.414898,-0.9000);
\fill[schur5] (1.439292,-1.0800) rectangle (1.512477,-0.9000);
\fill[schur5] (1.536872,-1.0800) rectangle (1.561266,-0.9000);
\fill[schur5] (1.610056,-1.0800) rectangle (1.658845,-0.9000);
\fill[schur5] (1.707635,-1.0800) rectangle (1.829609,-0.9000);
\fill[schur5] (1.854004,-1.0800) rectangle (1.902793,-0.9000);
\fill[schur5] (1.927188,-1.0800) rectangle (1.951583,-0.9000);
\fill[schur5] (2.000372,-1.0800) rectangle (2.073557,-0.9000);
\fill[schur5] (2.097952,-1.0800) rectangle (2.122346,-0.9000);
\fill[schur5] (2.195531,-1.0800) rectangle (2.219926,-0.9000);
\fill[schur5] (2.341899,-1.0800) rectangle (2.366294,-0.9000);
\fill[schur5] (2.488268,-1.0800) rectangle (2.512663,-0.9000);
\fill[schur5] (2.585847,-1.0800) rectangle (2.610242,-0.9000);
\fill[schur5] (4.195903,-1.0800) rectangle (4.220298,-0.9000);
\fill[schur5] (4.586220,-1.0800) rectangle (4.610615,-0.9000);
\fill[schur5] (4.756983,-1.0800) rectangle (4.805773,-0.9000);
\fill[schur5] (4.878957,-1.0800) rectangle (4.903352,-0.9000);
\fill[schur5] (5.074115,-1.0800) rectangle (5.098510,-0.9000);
\fill[schur5] (5.147300,-1.0800) rectangle (5.171695,-0.9000);
\fill[schur5] (5.732775,-1.0800) rectangle (5.757169,-0.9000);
\fill[schur5] (7.342831,-1.0800) rectangle (7.367225,-0.9000);
\fill[schur5] (7.928305,-1.0800) rectangle (7.952700,-0.9000);
\fill[schur5] (8.001490,-1.0800) rectangle (8.025885,-0.9000);
\fill[schur5] (8.196648,-1.0800) rectangle (8.221043,-0.9000);
\fill[schur5] (8.294227,-1.0800) rectangle (8.343017,-0.9000);
\fill[schur5] (8.489385,-1.0800) rectangle (8.513780,-0.9000);
\fill[schur5] (8.879702,-1.0800) rectangle (8.904097,-0.9000);
\fill[schur5] (10.489758,-1.0800) rectangle (10.514153,-0.9000);
\fill[schur5] (10.587337,-1.0800) rectangle (10.611732,-0.9000);
\fill[schur5] (10.733706,-1.0800) rectangle (10.758101,-0.9000);
\fill[schur5] (10.880074,-1.0800) rectangle (10.904469,-0.9000);
\fill[schur5] (10.977654,-1.0800) rectangle (11.002048,-0.9000);
\fill[schur5] (11.026443,-1.0800) rectangle (11.099628,-0.9000);
\fill[schur5] (11.148417,-1.0800) rectangle (11.172812,-0.9000);
\fill[schur5] (11.197207,-1.0800) rectangle (11.245996,-0.9000);
\fill[schur5] (11.270391,-1.0800) rectangle (11.392365,-0.9000);
\fill[schur5] (11.441155,-1.0800) rectangle (11.489944,-0.9000);
\fill[schur5] (11.538734,-1.0800) rectangle (11.563128,-0.9000);
\fill[schur5] (11.587523,-1.0800) rectangle (11.660708,-0.9000);
\fill[schur5] (11.685102,-1.0800) rectangle (11.709497,-0.9000);
\fill[schur5] (11.733892,-1.0800) rectangle (11.758287,-0.9000);
\fill[schur5] (11.782682,-1.0800) rectangle (11.807076,-0.9000);
\fill[schur5] (11.831471,-1.0800) rectangle (11.855866,-0.9000);
\fill[schur5] (11.880261,-1.0800) rectangle (11.904655,-0.9000);
\fill[schur5] (11.929050,-1.0800) rectangle (11.953445,-0.9000);
\fill[schur5] (12.026629,-1.0800) rectangle (12.051024,-0.9000);
\fill[schur5] (12.172998,-1.0800) rectangle (12.197393,-0.9000);
\node[anchor=west] at (13.219999999999999,-0.9900) {76};
\node[anchor=east] at (-0.12,-1.2600) {$C_6$};
\fill[black!4] (0,-1.3500) rectangle (13.1,-1.1700);
\fill[schur6] (3.659218,-1.3500) rectangle (3.683613,-1.1700);
\fill[schur6] (4.244693,-1.3500) rectangle (4.269088,-1.1700);
\fill[schur6] (4.317877,-1.3500) rectangle (4.342272,-1.1700);
\fill[schur6] (4.439851,-1.3500) rectangle (4.488641,-1.1700);
\fill[schur6] (4.513035,-1.3500) rectangle (4.537430,-1.1700);
\fill[schur6] (4.610615,-1.3500) rectangle (4.659404,-1.1700);
\fill[schur6] (4.708194,-1.3500) rectangle (4.756983,-1.1700);
\fill[schur6] (4.805773,-1.3500) rectangle (4.878957,-1.1700);
\fill[schur6] (4.903352,-1.3500) rectangle (4.927747,-1.1700);
\fill[schur6] (4.976536,-1.3500) rectangle (5.074115,-1.1700);
\fill[schur6] (5.098510,-1.3500) rectangle (5.147300,-1.1700);
\fill[schur6] (5.171695,-1.3500) rectangle (5.732775,-1.1700);
\fill[schur6] (5.757169,-1.3500) rectangle (7.342831,-1.1700);
\fill[schur6] (7.367225,-1.3500) rectangle (7.928305,-1.1700);
\fill[schur6] (7.952700,-1.3500) rectangle (8.001490,-1.1700);
\fill[schur6] (8.025885,-1.3500) rectangle (8.123464,-1.1700);
\fill[schur6] (8.172253,-1.3500) rectangle (8.196648,-1.1700);
\fill[schur6] (8.221043,-1.3500) rectangle (8.294227,-1.1700);
\fill[schur6] (8.343017,-1.3500) rectangle (8.391806,-1.1700);
\fill[schur6] (8.440596,-1.3500) rectangle (8.489385,-1.1700);
\fill[schur6] (8.562570,-1.3500) rectangle (8.586965,-1.1700);
\fill[schur6] (8.611359,-1.3500) rectangle (8.660149,-1.1700);
\fill[schur6] (8.757728,-1.3500) rectangle (8.782123,-1.1700);
\fill[schur6] (8.830912,-1.3500) rectangle (8.855307,-1.1700);
\fill[schur6] (9.416387,-1.3500) rectangle (9.440782,-1.1700);
\node[anchor=west] at (13.219999999999999,-1.2600) {151};
\draw[black!55,densely dashed] (6.550000,.24) -- (6.550000,-1.43);
\draw[black!50] (0.012197,-1.48) -- (0.012197,-1.54);
\node[anchor=north] at (0.012197,-1.55) {1};
\draw[black!50] (2.427281,-1.48) -- (2.427281,-1.54);
\node[anchor=north] at (2.427281,-1.55) {100};
\draw[black!50] (4.866760,-1.48) -- (4.866760,-1.54);
\node[anchor=north] at (4.866760,-1.55) {200};
\draw[black!50] (7.306238,-1.48) -- (7.306238,-1.54);
\node[anchor=north] at (7.306238,-1.55) {300};
\draw[black!50] (9.745717,-1.48) -- (9.745717,-1.54);
\node[anchor=north] at (9.745717,-1.55) {400};
\draw[black!50] (12.185196,-1.48) -- (12.185196,-1.54);
\node[anchor=north] at (12.185196,-1.55) {500};
\draw[black!50] (13.087803,-1.48) -- (13.087803,-1.54);
\node[anchor=north] at (13.087803,-1.55) {537};
\node[anchor=west] at (13.219999999999999,-1.73) {$x$};
\draw[black!60] (5.757169,-1.93) -- (5.757169,-2.00) -- (7.342831,-2.00) -- (7.342831,-1.93);
\node[anchor=north] at (6.550000,-2.02) {$[237,301]\subset C_6$};
\node[anchor=west,font=\small] at (-0.65,-2.70) {(b) Allowed colors in the folded frame, $f=\min(x,538-x)$};
\node[anchor=east] at (-0.12,-3.0400) {$C_1$};
\fill[black!4] (0,-3.1300) rectangle (13.1,-2.9500);
\fill[schur1!35] (0.000000,-3.1300) rectangle (1.801859,-2.9500);
\fill[schur1!35] (1.850558,-3.1300) rectangle (7.304833,-2.9500);
\fill[schur1!35] (7.353532,-3.1300) rectangle (11.444238,-2.9500);
\fill[schur1] (0.000000,-3.1300) rectangle (0.048699,-2.9500);
\node[anchor=east] at (-0.12,-3.3100) {$C_2$};
\fill[black!4] (0,-3.4000) rectangle (13.1,-3.2200);
\fill[schur2!35] (0.048699,-3.4000) rectangle (1.801859,-3.2200);
\fill[schur2!35] (1.850558,-3.4000) rectangle (7.304833,-3.2200);
\fill[schur2!35] (7.353532,-3.4000) rectangle (11.444238,-3.2200);
\node[anchor=east] at (-0.12,-3.5800) {$C_3$};
\fill[black!4] (0,-3.6700) rectangle (13.1,-3.4900);
\fill[schur3!35] (0.097398,-3.6700) rectangle (1.801859,-3.4900);
\fill[schur3!35] (1.850558,-3.6700) rectangle (7.304833,-3.4900);
\fill[schur3!35] (7.353532,-3.6700) rectangle (11.444238,-3.4900);
\node[anchor=east] at (-0.12,-3.8500) {$C_4$};
\fill[black!4] (0,-3.9400) rectangle (13.1,-3.7600);
\fill[schur4!35] (0.146097,-3.9400) rectangle (1.801859,-3.7600);
\fill[schur4!35] (1.850558,-3.9400) rectangle (7.304833,-3.7600);
\fill[schur4!35] (7.353532,-3.9400) rectangle (11.444238,-3.7600);
\node[anchor=east] at (-0.12,-4.1200) {$C_5$};
\fill[black!4] (0,-4.2100) rectangle (13.1,-4.0300);
\fill[schur5!35] (1.801859,-4.2100) rectangle (5.502974,-4.0300);
\fill[schur5!35] (7.353532,-4.2100) rectangle (11.492937,-4.0300);
\fill[schur5] (1.801859,-4.2100) rectangle (1.850558,-4.0300);
\fill[schur5] (11.444238,-4.2100) rectangle (11.492937,-4.0300);
\node[anchor=east] at (-0.12,-4.3900) {$C_6$};
\fill[black!4] (0,-4.4800) rectangle (13.1,-4.3000);
\fill[schur6!35] (7.304833,-4.4800) rectangle (11.444238,-4.3000);
\fill[schur6!35] (11.492937,-4.4800) rectangle (13.100000,-4.3000);
\fill[schur6] (7.304833,-4.4800) rectangle (7.353532,-4.3000);
\fill[schur6] (11.492937,-4.4800) rectangle (13.100000,-4.3000);
\draw[black!50] (0.024349,-4.61) -- (0.024349,-4.67);
\node[anchor=north] at (0.024349,-4.68) {1};
\draw[black!50] (1.826208,-4.61) -- (1.826208,-4.67);
\node[anchor=north] at (1.826208,-4.68) {38};
\draw[black!50] (5.478625,-4.61) -- (5.478625,-4.67);
\node[anchor=north] at (5.478625,-4.68) {113};
\draw[black!50] (7.329182,-4.61) -- (7.329182,-4.67);
\node[anchor=north] at (7.329182,-4.68) {151};
\draw[black!50] (11.468587,-4.61) -- (11.468587,-4.67);
\node[anchor=north] at (11.468587,-4.68) {236};
\draw[black!50] (13.075651,-4.61) -- (13.075651,-4.67);
\node[anchor=north] at (13.075651,-4.68) {269};
\node[anchor=west] at (13.219999999999999,-4.83) {$f$};
\end{tikzpicture}

\caption{The 537 certificate and its search frame. (a) A colored cell in row
$C_j$ marks an integer in that class; the right margin gives the class size.
The dashed line marks the center of reflection. (b) Light cells mark permitted
colors before propagation; dark cells mark forced colors. The row layout
follows the certificate diagrams of Heule~\cite[Fig.~1]{heule}; the entries
here are generated from our certificate and frame rules.}
\label{fig:structure}
\end{figure}

The first class of the example has a simple description as well. Put
\begin{align*}
A&=\{1,5,12,15,19,21,23,25,29,32,39,43,45,49\},\\
B&=A\cup(151-A)\cup\{151+a: a\in A,\ a\le 29\}\cup\{95\},
\end{align*}
a set of $38$ numbers, and $C^*=\{x\in\{1,\dots,537\}:\fold(x)\in B\}$, a set of $76$
numbers. The first class of the example is exactly $C^*$.

\begin{proposition}\label{prop:families}
The $131$ partitions of $\{1,\dots,537\}$ into six sum-free sets listed in
the companion file \texttt{solutions537.txt} are pairwise distinct, every class of each of them is
sum-free modulo $538$, and no one of them is obtained from another by a map
$x\mapsto rx \bmod 538$ with $\gcd(r,538)=1$ followed by renaming the colours.
They form the five families summarized in Table~\ref{tab:families}.
The example of Theorem~\ref{thm:537} is $\mathcal F_1.001$.
Each partition in $\mathcal F_1$--$\mathcal F_4$ agrees with the example in at least
$462$ of the $537$ positions.
\end{proposition}

\begin{proof}
All statements are finite checks, carried out by the companion script \texttt{check.py}.
\end{proof}

For Table~\ref{tab:families}, $d_{\rm pal}$ counts pairs $\{x,538-x\}$
with different colors, $\Delta_1=|C_1\mathbin{\triangle}C^*|$, and $h$ counts
folded positions $f>236$ outside color 6. Family $\mathcal F_5$ is measured after
$x\mapsto43x\bmod538$ and color renumbering; its original class minima are
$1,2,3,4,7,16$. Families $\mathcal F_1$--$\mathcal F_3$ satisfy frame $g$.
Family $\mathcal F_4$ satisfies the frame with the solid-middle condition removed.
These families describe the retained sample, not all possible 537-partitions.

\begin{table}[tbp]
\caption{The 131 retained partitions. Intervals give the observed minimum and
maximum within a family. Features of $\mathcal F_5$ use the normalization stated
in the text. All entries are computed from \texttt{solutions537.txt}.}
\label{tab:families}
\centering
\begin{tabular}{lrrrr}
\toprule
Family & Count & $d_{\rm pal}$ & $\Delta_1$ & $h$\\
\midrule
$\mathcal F_1$ & 107 & 0 & 0 & 0\\
$\mathcal F_2$ & 14  & 0 & 2 & 0\\
$\mathcal F_3$ & 5   & 2--12 & 0 & 0\\
$\mathcal F_4$ & 4   & 0 & 2--8 & 1--13\\
$\mathcal F_5$ & 1   & 5 & 3 & 5\\
\midrule
Total & 131 & & &\\
\bottomrule
\end{tabular}
\end{table}

\section{How to reproduce the search}\label{sec:method}
\paragraph{Encoding.} For a palindromic colouring use variables $v_{y,k}$ for
$y\in\{1,\dots,269\}$ and $k\in\{1,\dots,6\}$, meaning ``$\fold(x)=y$ has colour $k$''.
The clauses are: $v_{y,1}\vee\dots\vee v_{y,6}$ for each $y$; the clause
$\neg v_{\fold(x),k}\vee\neg v_{\fold(y),k}\vee\neg v_{\fold(x+y),k}$ for all
$1\le x\le y$, $x+y\le537$ and all $k$; and one unit clause for each condition of
Definition~\ref{def:frameg}. (A palindromic colouring with sum-free classes is
automatically sum-free modulo $538$.) The script \texttt{gen.py} in
the companion repository writes this CNF and decodes a model.

\paragraph{Decoder caveat.}
The supplied historical generator does not impose at-most-one-colour clauses.
Its decoder selects the lowest-numbered true colour at each position. This
preserves sum-freeness, but can discard a colour required by a positive unit
clause, so the output need not satisfy frame $g$. For example, a satisfying
model can decode positions $237$ and $301$ as colour $1$, although the frame
requires colour $6$ there. This does not affect the independently checked
certificates reported here. To guarantee the frame after decoding, add
$\neg v_{y,i}\vee\neg v_{y,j}$ for each position $y$ and $1\le i<j\le6$,
or choose a required colour first and check all frame conditions after decoding.
The repository includes a script that reproduces this caveat.

\paragraph{Solving.} With Kissat~4.0.4~\cite{kissat}:
\begin{lstlisting}
python3 gen.py cnf > g.cnf
kissat --seed=1 g.cnf > g.out
python3 gen.py decode g.out
\end{lstlisting}
On one core this took $286$\,s with seed~1 and $4$\,s with seed~2; both solutions lie
in $\mathcal F_1$, and different seeds give different members.

\paragraph{The other families.} For $\mathcal F_3$ drop palindromy: use one variable
per $x\in\{1,\dots,537\}$, forbid monochromatic $x+y\equiv z\pmod{538}$, and impose
the frame conditions on $\fold(x)$ (\texttt{gen.py --nopal}; $41$\,s with seed~2,
while seed~1 did not finish in $15$ minutes). For $\mathcal F_4$ drop the condition that
colour~6 fills $\fold(x)>236$ (\texttt{gen.py --nomid}); this allows $\mathcal F_4$,
but a given run may still return a member of $\mathcal F_1$.
$\mathcal F_5$ came from a list-colouring run: let $c$ be the example and
$c_m(x)=c(mx\bmod 538)$ for $m=3,25,221$, with colours renamed to agree with $c$ as
much as possible; each $x$ may only take one of the colours $c_3(x),c_{25}(x),
c_{221}(x)$, and no other condition is imposed.

\paragraph{Where frame $g$ came from.} An evolutionary search over depth-first
search programs, in the spirit of FunSearch~\cite{funsearch}, produced valid
colourings of $\{1,\dots,n\}$ for $n$ up to $528$ from scratch, without seed
colourings. Those programs restrict the search by rules that ban a colour below a
given value of $\fold$, and the rules used by the best colourings had the shape of
Definition~\ref{def:frameg}, with $Q$ and $d$ as parameters. At $n=537$ a sweep over
$140\le Q\le160$ found no solution for $39\le d\le45$ and the first one at
$(Q,d)=(151,38)$.

The work was done by a group of AI agents (Claude Code sessions, with OpenAI Codex
for second opinions) on a $192$-core server over about a day and a half; every
result was recorded with a command that reproduces it and checked by a second agent.

\paragraph{Data and code.}
The companion repository contains the full list of $131$ partitions, a separate
PDF table, the example certificate, the $5372$-element eight-colouring,
\texttt{check.py}, \texttt{gen.py}, and \texttt{paper/make\_figures.py}.
The figures are vector graphics generated from the certificates. To check the results, run
\texttt{python3 check.py} from the repository root with Python~3 and its
standard library. The paper retains one complete partition; the remaining
partitions and code are supplementary material.

\begin{thebibliography}{99}
\small\raggedright
\bibitem{schur} I.~Schur, \"Uber die Kongruenz $x^m+y^m\equiv z^m \pmod p$,
\emph{Jahresber.\ Deutsch.\ Math.-Verein.} \textbf{25} (1916), 114--117.
\bibitem{abbott} H.~L.~Abbott and D.~Hanson, A problem of Schur and its
generalizations, \emph{Acta Arith.} \textbf{20} (1972), 175--187.
\bibitem{rowley} F.~Rowley, An improved lower bound for $S(7)$ and some interesting
templates, arXiv:2107.03560 (2021).
\bibitem{ageron} R.~Ageron, P.~Casteras, T.~Pellerin, Y.~Portella, A.~Rimmel and
J.~Tomasik, New lower bounds for Schur and weak Schur numbers, arXiv:2112.03175
(2021).
\bibitem{bengone} N.~Bengone, A.~Brouk, M.~Grinsztajn, T.~Helbert, B.~Lugherini,
A.~Rimmel and J.~Tomasik, Shifted S-templates and improved lower bounds for Schur
numbers, arXiv:2607.15034 (2026).
\bibitem{xxer} X.~Xu, Z.~Xie, G.~Exoo and S.~P.~Radziszowski, Constructive lower
bounds on classical multicolor Ramsey numbers, \emph{Electron.\ J.\ Combin.}
\textbf{11} (2004), \#R35.
\bibitem{fs} H.~Fredricksen and M.~M.~Sweet, Symmetric sum-free partitions and lower
bounds for Schur numbers, \emph{Electron.\ J.\ Combin.} \textbf{7} (2000), \#R32. \href{https://doi.org/10.37236/1510}{doi:10.37236/1510}.
\bibitem{heule} M.~J.~H.~Heule, Schur number five, \emph{Proc.\ AAAI} \textbf{32}
(2018). \href{https://doi.org/10.1609/aaai.v32i1.12209}{doi:10.1609/aaai.v32i1.12209}.
\bibitem{kissat} A.~Biere, K.~Fazekas, M.~Fleury and M.~Heisinger, CaDiCaL, Kissat,
Paracooba, Plingeling and Treengeling entering the SAT Competition 2020,
\emph{Proc.\ SAT Competition 2020}.
\bibitem{funsearch} B.~Romera-Paredes et al., Mathematical discoveries from program
search with large language models, \emph{Nature} \textbf{625} (2024).
\end{thebibliography}
\end{document}
